%! Solving by Square Roots & Completing the Square — Unit 3, Lesson 3.3 — Group Activity (Answer Key)
\documentclass[10pt]{article}
\usepackage{algebra2-article}
\usepackage{algebra2-key}   % pulls in -boxes; do NOT also load -boxes
\newcommand{\TallMath}[1]{$\displaystyle #1\rule[-1.4em]{0pt}{3.2em}$}
\newcommand{\lgrid}[4]{%
  \draw[step=1, linegray!40, very thin] (#1,#3) grid (#2,#4);
  \draw[->, semithick, charcoal] (#1,0)--(#2,0) node[below right, font=\scriptsize]{$x$};
  \draw[->, semithick, charcoal] (0,#3)--(0,#4) node[left, font=\scriptsize]{$y$};}

\begin{document}

\pageheader{Unit 3, Lesson 3.3}{Group Activity --- Answer Key}
\namepartnerperiod

\vspace{0.08in}

\begin{headlinebox}{forestbg}
\small\textbf{Undo the Square.} When a quadratic won't factor, you can still find its roots: isolate a
square and take $\pm\sqrt{\ }$, or \emph{complete the square} to build one. With your partner, solve, keep
\emph{both} roots in simplest radical form, and check each answer against the graph.
\end{headlinebox}

\vspace{0.06in}

\begin{tcolorbox}[colback=white, colframe=black!40, title=\textbf{Tier R --- Remediate},
    fonttitle=\bfseries, arc=1.5mm, left=3mm, right=3mm, top=2mm, bottom=2mm, breakable]
\textbf{1. Square Root Property.} Solve for \emph{both} roots.
\[
  x^2=36:\ x=\textcolor{keyred}{\mathbf{\pm6}} \qquad x^2=100:\ x=\textcolor{keyred}{\mathbf{\pm10}} \qquad x^2=1:\ x=\textcolor{keyred}{\mathbf{\pm1}}
\]

\textbf{2. Squared binomials.} Root both sides, then undo the shift.
\[
  (x-1)^2=25:\ x=\textcolor{keyred}{\mathbf{6,\,-4}} \qquad (x+3)^2=4:\ x=\textcolor{keyred}{\mathbf{-1,\,-5}}
\]

\textbf{3. Match to a graph.} Solve each equation, then write the letter of the parabola whose
$x$-intercepts match.
\begin{center}
\begin{tikzpicture}[scale=0.36]
  % A: roots -2, 2  (y = x^2 - 4)
  \begin{scope}
    \lgrid{-4}{4}{-5}{3}
    \begin{scope}\clip (-4,-5) rectangle (4,3);
      \draw[very thick, forest, domain=-2.6:2.6, samples=60, smooth] plot (\x,{\x*\x-4});
    \end{scope}
    \fill[charcoal] (-2,0) circle (2.5pt); \fill[charcoal] (2,0) circle (2.5pt);
    \node[charcoal, font=\scriptsize] at (0,-6.0){\textbf{A}};
  \end{scope}
  % B: roots -3, 3  (y = x^2 - 9)
  \begin{scope}[xshift=10cm]
    \lgrid{-4}{4}{-10}{3}
    \begin{scope}\clip (-4,-10) rectangle (4,3);
      \draw[very thick, forest, domain=-3.4:3.4, samples=70, smooth] plot (\x,{\x*\x-9});
    \end{scope}
    \fill[charcoal] (-3,0) circle (2.5pt); \fill[charcoal] (3,0) circle (2.5pt);
    \node[charcoal, font=\scriptsize] at (0,-11.5){\textbf{B}};
  \end{scope}
  % C: roots -5, 1  (y = (x+2)^2 - 9)
  \begin{scope}[xshift=20cm]
    \lgrid{-6}{2}{-10}{3}
    \begin{scope}\clip (-6,-10) rectangle (2,3);
      \draw[very thick, forest, domain=-5.4:1.4, samples=70, smooth] plot (\x,{(\x+2)*(\x+2)-9});
    \end{scope}
    \fill[charcoal] (-5,0) circle (2.5pt); \fill[charcoal] (1,0) circle (2.5pt);
    \node[charcoal, font=\scriptsize] at (-2,-11.5){\textbf{C}};
  \end{scope}
\end{tikzpicture}
\end{center}
\renewcommand{\arraystretch}{1.5}
\begin{tabularx}{\linewidth}{@{}l X c@{}}
\textbf{Equation} & \textbf{Roots} & \textbf{Graph} \\
\hline
$x^2=9$       & \ans{$\pm3$}   & \ans{B} \\
$(x+2)^2=9$   & \ans{$1,\ -5$} & \ans{C} \\
$x^2=4$       & \ans{$\pm2$}   & \ans{A} \\
\end{tabularx}
\end{tcolorbox}

\begin{tcolorbox}[colback=white, colframe=black!40, title=\textbf{Tier A --- Approaching Proficiency},
    fonttitle=\bfseries, arc=1.5mm, left=3mm, right=3mm, top=2mm, bottom=2mm, breakable]
\textbf{1. Isolate, then root.} Simplest radical form; keep both roots.
\[
  3x^2-24=0:\ x=\textcolor{keyred}{\mathbf{\pm2\sqrt2}} \qquad (x-4)^2=20:\ x=\textcolor{keyred}{\mathbf{4\pm2\sqrt5}}
\]
\[
  x^2+5=41:\ x=\textcolor{keyred}{\mathbf{\pm6}} \qquad 2x^2=54:\ x=\textcolor{keyred}{\mathbf{\pm3\sqrt3}}
\]

\textbf{2. Complete the square.} Show the completed square, then the roots.
\[
  x^2+8x+3=0:\ (x+\textcolor{keyred}{\mathbf{4}})^2=\textcolor{keyred}{\mathbf{13}},\ \ x=\textcolor{keyred}{\mathbf{-4\pm\sqrt{13}}}
\]
\[
  x^2-6x-2=0:\ (x-\textcolor{keyred}{\mathbf{3}})^2=\textcolor{keyred}{\mathbf{11}},\ \ x=\textcolor{keyred}{\mathbf{3\pm\sqrt{11}}}
\]

\textbf{3. Vertex form.} Complete the square to rewrite each and name the vertex.
\begin{center}
\renewcommand{\arraystretch}{1.6}
\begin{tabularx}{\linewidth}{@{}l X X@{}}
\textbf{Standard form} & \textbf{Vertex form $(x-h)^2+k$} & \textbf{Vertex $(h,k)$} \\
\hline
$x^2+8x+10$  & \ans{$(x+4)^2-6$}  & \ans{$(-4,-6)$} \\
$x^2-10x+18$ & \ans{$(x-5)^2-7$}  & \ans{$(5,-7)$}  \\
\end{tabularx}
\end{center}
\end{tcolorbox}

\begin{tcolorbox}[colback=white, colframe=black!40, title=\textbf{Tier E --- Extension},
    fonttitle=\bfseries, arc=1.5mm, left=3mm, right=3mm, top=2mm, bottom=2mm, breakable]
\textbf{1. Leading coefficient $a\ne1$ (divide first).} Complete the square and solve.
\[
  2x^2-4x-3=0:\ x=\textcolor{keyred}{\mathbf{1\pm\tfrac{\sqrt{10}}{2}}}
\]
\emph{Work:} divide by $2$: $x^2-2x-\tfrac32=0 \Rightarrow x^2-2x=\tfrac32 \Rightarrow (x-1)^2=\tfrac52
\Rightarrow x-1=\pm\tfrac{\sqrt{10}}{2}$.

\textbf{2. Choose the smartest method.} For each, name the method (\emph{factoring / square roots /
completing the square}), then solve. \emph{Justify} your choice.
\begin{center}
\renewcommand{\arraystretch}{1.6}
\begin{tabularx}{\linewidth}{@{}l X X@{}}
\textbf{Equation} & \textbf{Method (why?)} & \textbf{Roots} \\
\hline
$x^2-49=0$      & \ans{square roots (no $bx$ term)} & \ans{$\pm7$} \\
$x^2-7x+12=0$   & \ans{factoring (factors nicely)}  & \ans{$3,\ 4$} \\
$x^2+2x-5=0$    & \ans{completing the square (won't factor)} & \ans{$-1\pm\sqrt6$} \\
\end{tabularx}
\end{center}

\vspace{4pt}
\textbf{3. Model it.} A stone is dropped from a $100$-ft ledge. Its height (ft) after $t$ seconds is
$h(t)=-16t^2+100$.
\begin{enumerate}[label=(\alph*), leftmargin=1.9em, itemsep=2pt]
  \item How high is the stone at $t=1$ second? \ans{$-16+100=84$ ft}
  \item When does it hit the ground? Solve $-16t^2+100=0$ by the square-root method. \ansline{$16t^2=100 \Rightarrow t^2=6.25 \Rightarrow t=\pm2.5$; the stone lands at $t=2.5$ s.}
  \item One solution is impossible here. Which, and why? \ansline{Reject $t=-2.5$: time cannot be negative.}
\end{enumerate}
\end{tcolorbox}

\begin{teachernote}
Tier~R keeps radicands to perfect squares so the focus is the $\pm$ and undoing a shift (item~2), plus the
graph link (item~3: $x^2=9\to$B, $(x+2)^2=9\to$C with roots $1,-5$, $x^2=4\to$A). Tier~A item~1 is
``isolate first, then simplify'' ($3x^2-24=0\Rightarrow x^2=8\Rightarrow\pm2\sqrt2$); watch students who
root before dividing. Item~3 is the vertex-form payoff --- verify by expanding. In Tier~E, the $a\ne1$ case
must \emph{divide by $a$ first}; the choose-a-method table rewards recognizing structure (no $bx\Rightarrow$
roots; nice factors $\Rightarrow$ factor; otherwise complete the square). In the model, insist on rejecting
$t=-2.5$ \emph{with a reason}.
\end{teachernote}

\end{document}
